\section{Conclusion}\label{sec:conclusion}

Intergenerational persistence in the United States rises steeply with the age at which children are observed, and
the rise is concentrated among White sons. The White--Black gap in the elasticity of sons' earnings with respect to
parental income is zero in the late twenties and about 0.3 in the forties; with one year of parental income it opens
by about 0.09 log points in four public data sets, 0.19 corrected for measurement error, and it survives attrition weights,
period effects and a dollar-weighted elasticity. In ranks of family income the racial gap is a level gap that is present from the
start; in logs there is in addition a gap in persistence that is built over the career. It is a gap in wages, and it
is built in the first decade of the career: in the jobs sons hold at 25--34, to which access is unequal by background
and by race, and in wage growth within them, which runs through a return to schooling that rises with age for White
sons and much less for Black sons. Sons of richer White parents also converge on their fathers' occupations, and the
convergence on the father's exact job is the inheritance of a family business.

Shocks and policies at the start of the career do not create the gap. Recessions at entry leave average access
to good jobs unchanged and the gap opens among sons who entered tight markets too, and public employment, where Black
sons start more often, carries none of the divergence. A childhood shock re-scored by the son's age---the father's
displacement---shows a measured effect that does not change between 25--34 and 35--44, because it acts on the level of
earnings, not on their growth.

Two implications follow. For measurement, an IGE measured before the mid-thirties understates persistence, and
understates it more for White than for Black families, so the racial comparison depends on the age of measurement in a
way the level of persistence alone does not reveal. For policy, interventions that act on the starting point of the
career will look more effective at 30 than at 45. The life-cycle model of Section~\ref{sec:model} says why: among four
explanations it selects the one in which parental income raises the growth of human capital for White sons and not for
Black sons, together with a pay-off of schooling that is a wage level for Black sons and a growth rate for White sons,
and it rejects a racial difference in offers of good jobs: the lower Black share in good jobs at equal parental
income is reproduced without one. Equalising schooling, access or entry conditions closes less
than a tenth of the gap at 45; what closes it is the pay-off of background and schooling in growth, which the next
step of this work has to explain.
